\documentclass{exam}
\usepackage{../../shah350}

\hypersetup{colorlinks=true, linktoc=section, linkcolor=blue}

\pagestyle{headandfoot}
\firstpageheadrule
\runningheadrule
\firstpageheader{Prof. Lepowsky \\ Linear Algebra}{Homework \#4}{Jeevan Shah}
\runningheader{Linear Algebra \\ Homework \#4}{}{Shah}
\firstpagefooter{}{\thepage}{}
\runningfooter{ }{\thepage}{ }

\printanswers


\begin{document}
  \colorbox{red}{\underline{$1.4.15$}:}
  \begin{solution}
    \begin{proof}
      Consider $S_1, S_2$ subsets of a vector space $V$ and let $\vec{x} \in \Span{S_1 \cap S_2}$. If $S_1 \cap S_2 = \varnothing$ then $\vec{x} = \vec{0}$, which lies in every span. Otherwise, by definition
      \[
        \vec{x} = \sum_{i=1}^{n} a_i s_i
      \]
      for some $s_1, \dots, s_n \in S_1 \cap S_2$ and $a_i \in \FF$. By definition of intersection each $s_i \in S_1$, so $\vec{x}$ is a linear combination of vectors in $S_1$ and hence $\vec{x} \in \Span{S_1}$. An identical argument shows that $\vec{x} \in \Span{S_2}$. Thus $\Span{S_1 \cap S_2} \subseteq \Span{S_1} \cap \Span{S_2}$ as was to be shown.
    \end{proof}

    If $S_1 = \varnothing$ then $\Span{S_1 \cap S_2} = \Span{\varnothing} = \braces{\vec{0}} = \Span{\varnothing} \cap \Span{S_2}$ for any $S_2$. Now to see an example in which they are not equal, take $V = \RR^2, S_1 = \braces{(1, 0)}$ and $S_2 = \braces{(2, 0)}$ so that 
    \[
      S_1 \cap S_2 = \varnothing \Rightarrow \Span{S_1 \cap S_2} = \braces{\vec{0}} 
    \]
    but 
    \[
      \Span{S_1} = \Span{S_2} = \Span{S_1} \cap \Span{S_2} = \braces{x\begin{pmatrix}
        1 \\ 0
      \end{pmatrix} : x \in \RR} \neq \braces{\vec{0}}.
    \]
  \end{solution}

  \colorbox{red}{\underline{$1.5.15$}:}
  \begin{solution}
    \begin{proof}
      $(\Rightarrow)$ First, assume $S$ is linearly dependent. Then there exist scalars $a_1, \dots, a_n$, not all zero, such that $a_1u_1 + \dots + a_nu_n = \vec{0}$. Let $j$ be the largest index with $a_j \neq 0$. If $j = 1$ then $a_1u_1 = \vec{0}$ with $a_1 \neq 0$, so $u_1 = \vec{0}$. Otherwise $j \geq 2$ and
      \[
        u_j = -a_j^{-1}\left(a_1u_1 + \dots + a_{j-1}u_{j-1}\right),
      \]
      so $u_{k+1} \in \Span{u_1, \dots, u_k}$ for $k = j - 1$. \\
      $(\Leftarrow)$ Conversely, assume $u_1 = \vec{0}$ or $u_{k+1} \in \Span{u_1, \dots, u_k}$ for some $k$. If $u_1 = \vec{0}$ then $S$ is linearly dependent since $1 \cdot u_1 = \vec{0}$ is a non-trivial representation of $\vec{0}$. If $u_{k+1} \in \Span{u_1, \dots, u_k}$ then $u_{k+1} = a_1u_1 + \dots + a_ku_k$ for some scalars $a_i$, so $a_1u_1 + \dots + a_ku_k + (-1)u_{k+1} = \vec{0}$ is a non-trivial representation of $\vec{0}$ since the coefficient of $u_{k+1}$ is $-1 \neq 0$. Hence $S$ is linearly dependent.
    \end{proof}

    $k$ need not be $0$: the case $k = 0$ would say $u_1 \in \Span{\varnothing} = \braces{\vec{0}}$, which is exactly the case $u_1 = \vec{0}$.
  \end{solution}

  \colorbox{red}{\underline{$2.1.16$}:}
  \begin{solution}
    \begin{proof}
      In order to show that $T$ is onto we must show that for all $f \in P(\RR)$ there exists $g \in P(\RR)$ such that $T(g) = f$. However, since $T$ is defined by $T(f(x)) = f'(x)$ we can simply apply the fundamental theorem of calculus to take 
      \[
        g(x) = \int_{0}^{x} f(t)\,\mathrm{d}t
      \]
      so that 
      \[
        T(g(x)) = T\left(\int_{0}^{x}f(t)\,\mathrm{d}t\right) = \left(\int_{0}^{x}f(t)\,\mathrm{d}t\right)' = f(x).
      \]
      Now, to show that $T$ is not one-to-one consider $h(x) = 3x + 1$ and $k(x) = 3x + 2$ so that 
      \[
        T(h(x)) = 3 = T(k(x)) 
      \]
      but $h(x) \neq k(x)$ and so $T$ cannot be one-to-one.
    \end{proof}
  \end{solution}

  \colorbox{red}{\underline{$2.2.2$}:}
  \begin{solution}
    \begin{parts}
      \part We have that 
      \[
        T(e_1) = \begin{pmatrix}
          2 \\ 3 \\ 1
        \end{pmatrix}
        \quad\text{and}\quad
        T(e_2) = \begin{pmatrix}
          -1 \\ 4 \\ 0
        \end{pmatrix}
      \]
      so 
      \[
        \left[T\right]_{\beta}^{\gamma} = \begin{pmatrix}
          2 & -1 \\
          3 & 4 \\
          1 & 0 
        \end{pmatrix} = A
      \]
      and so the matrix representation of $T$ in the standard ordered bases for $\RR^{2}$ and $\RR^{3}$ is the same as the matrix representation of $T$ in the form $T=L_A$.

      \part We have that 
      \[
        T(e_1) = \begin{pmatrix}
          2 \\ 1
        \end{pmatrix}
        \quad\text{and}\quad
        T(e_2) = \begin{pmatrix}
          3 \\ 0
        \end{pmatrix}
        \quad\text{and}\quad
        T(e_3) = \begin{pmatrix}
          -1 \\ 1
        \end{pmatrix}
      \]
      so 
      \[
        \left[T\right]_{\beta}^{\gamma} = \begin{pmatrix}
          2 & 3 & -1 \\
          1 & 0 & 1
        \end{pmatrix} = A
      \]
      and so the matrix representation of $T$ in the standard ordered bases for $\RR^{3}$ and $\RR^{2}$ is the same as the matrix representation of $T$ in the form $T=L_A$.

      \part We have that 
      \[
        T(e_1) = 2
        \quad\text{and}\quad
        T(e_2) = 1
        \quad\text{and}\quad
        T(e_3) = -3
      \]
      so 
      \[
        \left[T\right]_{\beta}^{\gamma} = \begin{pmatrix}
          2 & 1 & -3
        \end{pmatrix} = A
      \]
      and so the matrix representation of $T$ in the standard ordered bases for $\RR^{3}$ and $\RR$ is the same as the matrix representation of $T$ in the form $T=L_A$.
    \end{parts}
  \end{solution}


  \colorbox{red}{\underline{$2.2.4$}:}
  \begin{solution}
    We have that 
    \[
      T(e_1) = 1 
      \quad\text{and}\quad
      T(e_2) = 1+x^2
      \quad\text{and}\quad 
      T(e_3) = 0
      \quad\text{and}\quad 
      T(e_4) = 2x
    \]
    so 
    \[
      \left[T\right]_{\beta}^{\gamma} = \begin{pmatrix}
        1 & 1 & 0 & 0 \\
        0 & 0 & 0 & 2 \\
        0 & 1 & 0 & 0 
      \end{pmatrix}
    \]
  \end{solution}

  \newpage

  \colorbox{red}{\underline{$2.2.5d$}:}
  \begin{solution}
    We have that 
    \[
      T(e_1) = 1 
      \quad\text{and}\quad 
      T(e_2) = 2 
      \quad\text{and}\quad
      T(e_3) = 4  
    \]
    so 
    \(
      \left[T\right]_{\beta}^{\gamma} = \begin{pmatrix}
        1 & 2 & 4
      \end{pmatrix}
    \)
  \end{solution}

  \colorbox{red}{\underline{$2.3.3$}:}
  \begin{solution}
    \begin{parts}
      \part Let $e_1, e_2$, and $e_3$ be the standard basis vectors for $P_2(\RR)$. First, in order to find $[U]_{\beta}^{\gamma}$ and $[T]_\beta$ we see that 
      \[
        U(e_1) = \begin{pmatrix}
          1 \\ 0 \\ 1
        \end{pmatrix},
        \quad
        U(e_2) = \begin{pmatrix}
          1 \\ 0 \\ -1
        \end{pmatrix},
        \quad\text{and}\quad 
        U(e_3) = \begin{pmatrix}
          0 \\ 1 \\ 0
        \end{pmatrix}
      \] 
      and 
      \[
        T(e_1) = 2, \quad T(e_2) = 3 + 3x, \quad\text{and}\quad T(e_3) = 6x + 4x^2.
      \]
      So, 
      \[
        [U]_{\beta}^{\gamma} = \begin{pmatrix}
          1 & 1 & 0 \\
          0 & 0 & 1 \\
          1 & -1 & 0
        \end{pmatrix}
        \quad\text{and}\quad 
        [T]_\beta = \begin{pmatrix}
          2 & 3 & 0 \\ 
          0 & 3 & 6 \\
          0 & 0 & 4
        \end{pmatrix}
      \]
      and 
      \[
        UT(e_1) = \begin{pmatrix}
          2 \\ 0 \\ 2
        \end{pmatrix},
        \quad UT(e_2) = \begin{pmatrix}
          6 \\ 0 \\ 0
        \end{pmatrix},
        \quad\text{and}\quad 
        UT(e_3) = \begin{pmatrix}
          6 \\ 4 \\ -6
        \end{pmatrix}.
      \]
      Thus, 
      \[
        [UT]_\beta^\gamma = \begin{pmatrix}
          2 & 6 & 6 \\
          0 & 0 & 4 \\
          2 & 0 & -6
        \end{pmatrix}.
      \]
      Now, by Thm.~2.11 we have that 
      \[
        [UT]_\beta^\gamma = [U]_\beta^\gamma [T]_\beta = \begin{pmatrix}
          1 & 1 & 0 \\
          0 & 0 & 1 \\
          1 & -1 & 0
        \end{pmatrix}
        \begin{pmatrix}
          2 & 3 & 0 \\
          0 & 3 & 6 \\
          0 & 0 & 4
        \end{pmatrix}
        = \begin{pmatrix}
          2 & 6 & 6 \\
          0 & 0 & 4 \\
          2 & 0 & -6
        \end{pmatrix}
      \]
      which is exactly the matrix that we found earlier.

      \part Clearly we have 
      \[
        [h(x)]_\beta = \begin{pmatrix}
          3 \\ -2 \\ 1
        \end{pmatrix}
        \quad\text{and}\quad 
        [U(h(x))]_\gamma = \begin{pmatrix}
          1 \\ 1 \\ 5
        \end{pmatrix}.
      \]
      Using Thm.~2.14 we check:  
      \[
        [U(h(x))]_\gamma = [U]_\beta^\gamma [h(x)]_\beta = \begin{pmatrix}
          1 & 1 & 0 \\
          0 & 0 & 1 \\
          1 & -1 & 0 
        \end{pmatrix}
        \begin{pmatrix}
          3 \\ -2 \\ 1
        \end{pmatrix}
        = \begin{pmatrix}
          1 \\ 1 \\ 5
        \end{pmatrix}
      \]
    \end{parts}
  \end{solution}

  \newpage 

  \colorbox{red}{\underline{$2.3.4d$}:}
  \begin{solution}
    We first find $[T]_\beta^\gamma$ and $[f(x)]_\beta$:
    \[
      [T]_\beta^\gamma = \begin{pmatrix}
        1 & 2 & 4
      \end{pmatrix}
      \quad\text{and}\quad
      [f(x)]_\beta = \begin{pmatrix}
        6 \\ -1 \\ 2
      \end{pmatrix}
    \]
    thus
    \[
      [T(f(x))]_\gamma = [T]_\beta^\gamma [f(x)]_\beta = \begin{pmatrix}
        1 & 2 & 4
      \end{pmatrix}
      \begin{pmatrix}
        6 \\ -1 \\ 2
      \end{pmatrix}
      = \begin{pmatrix}
        12
      \end{pmatrix}.
    \]
  \end{solution}

  \colorbox{red}{\underline{$2.4.2$}:}
  \begin{solution}
    \begin{parts}
      \setcounter{partno}{2}
      \part We have that the matrix representation of $T$ is 
      \[
        A = \begin{pmatrix}
          3 & 0 & -2 \\
          0 & 1 & 0 \\
          3 & 4 & 0 
        \end{pmatrix} 
        \xrightarrow{r_3 \to r_3 - r_1}
        \begin{pmatrix}
          3 & 0 & -2 \\
          0 & 1 & 0 \\
          0 & 4 & 2
        \end{pmatrix}
        \xrightarrow{r_3 \to r_3 - 4r_2}
        \begin{pmatrix}
          3 & 0 & -2 \\
          0 & 1 & 0 \\
          0 & 0 & 2
        \end{pmatrix}
      \]
      and so clearly $\rank{A} = 3 = \dim{\RR^{3}}$ and thus $T$ is invertible by Thm.~2.5.

      \part As shown in ex(2.1.16), $T$ is not one-to-one and thus cannot be invertible.
    \end{parts}
  \end{solution}

  \colorbox{red}{\underline{$2.4.7$}:}
  \begin{solution}
    \begin{parts}
      \part \begin{proof}
        Suppose for the sake of contradiction that $A$ is invertible. Then, 
        \begin{align*}
          A^2 = 0 \Rightarrow AA = 0 &\Rightarrow A^{-1}(AA) = A^{-1}0 \\
          &\Rightarrow (A^{-1}A)A = 0 \\
          &\Rightarrow IA = 0 \\
          &\Rightarrow (IA)A^{-1} = 0A^{-1} \\
          &\Rightarrow I(AA^{-1}) = 0 \\
          &\Rightarrow II = 0 \\
          &\Rightarrow I = 0
        \end{align*}
        which is a contradiction and thus $A$ is not invertible.
      \end{proof}

      \part $A$ cannot be invertible, for otherwise
      \[
        AB = 0 \Rightarrow A^{-1}(AB) = A^{-1}0 \Rightarrow (A^{-1}A)B = 0 \Rightarrow IB = 0 \Rightarrow B = 0 
      \]
      which is a contradiction as $B$ is specifically nonzero and so $A$ cannot be invertible.
    \end{parts}
  \end{solution}

  \footnotetext{\LaTeX\, code for this document can be found on github \href{https://github.com/jeevanshah07/MATH350/blob/main/homework/homework4/main.tex}{\underline{here}}}
\end{document}